\documentclass[12pt,a4paper]{article}
\usepackage[utf8]{inputenc}
\usepackage[T1]{fontenc}
\usepackage[spanish]{babel}
\usepackage{amsmath}
\usepackage{amssymb}
\usepackage{geometry}
\usepackage{hyperref}
\usepackage{bookmark}
\usepackage{fancyhdr}
\usepackage{enumitem}
\usepackage{booktabs}
\usepackage{titlesec}
\geometry{margin=2.5cm}

% Encabezado estilo notas de clase
\setlength{\headheight}{15pt}
\pagestyle{fancy}
\fancyhf{}
\fancyhead[L]{Fundamentos de Valoración de Activos}
\fancyhead[R]{Problem Set \#2}
\fancyfoot[C]{\thepage}
\fancyfoot[L]{TA Session \#2}
\fancyfoot[R]{Jesus Villota Miranda}
\renewcommand{\footrulewidth}{0.4pt}

% Encabezado
\title{\textbf{Problem Set 2 Unificado}\\\large Fundamentos de Valoración de Activos}
\author{Jesus Villota Miranda\\CEMFI}
\date{\today}

\newcommand{\E}{\mathbb{E}}
\newcommand{\EUR}[1]{\ensuremath{#1~\text{EUR}}}

% Formato de ejercicios
\newenvironment{ejercicio}[1]{\subsection*{Ejercicio #1}}{}
\newenvironment{solucion}{\paragraph{Solución.}}{}

\begin{document}

% \maketitle

% \begin{center}
% \textbf{Objetivo.} Este problem set integra los mejores ejercicios de las cuatro versiones previas para cubrir: (i) valoración de acciones, (ii) valoración corporativa vía FCFF/FCFE, (iii) medición de rentabilidad y riesgo, y (iv) diversificación de carteras. Las dificultades están escalonadas desde recordatorios básicos hasta desafíos avanzados con extensiones conceptuales.
% \end{center}

% \tableofcontents
% \newpage

\section{Valoración de acciones}

\begin{ejercicio}{1 (Calentamiento: dividendos finitos)}
La acción de la empresa Lince pagará dividendos de $D_1=\,\EUR{2.20}$, $D_2=\,\EUR{2.38}$ y $D_3=\,\EUR{2.56}$. Al final del año 3 se espera venderla por $P_3=\,\EUR{34}$. Si la rentabilidad requerida es $r=11\%$, calcule el precio justo hoy $P_0$. ¿Qué porcentaje del valor proviene de los dividendos frente al precio de venta?
\end{ejercicio}

\begin{ejercicio}{2 (Crecimiento implícito y sensibilidad)}
Una acción vale hoy $P_0=\,\EUR{30}$ y pagará un dividendo próximo de $D_1=\,\EUR{1.50}$. La rentabilidad requerida es $r=9\%$.
\begin{enumerate}[label=\alph*)]
  \item Calcule la tasa de crecimiento constante $g$ implícita en la valoración de acuerdo con el modelo de Gordon.
  \item Si la empresa anuncia que el próximo dividendo ($D_1$) será un 5\% mayor del previsto inicialmente manteniendo $r$, ¿cuál sería el nuevo precio justo? Compare con $P_0$.
\end{enumerate}
\end{ejercicio}

\begin{ejercicio}{3 (Modelo de dos etapas y descomposición de rentabilidad)}
El dividendo actual es $D_0=\,\EUR{1.10}$. Se espera que los dividendos crezcan un 15\% anual durante tres años y, a partir de entonces, un ritmo constante $g_2$ de forma indefinida. La tasa de descuento requerida es $r=10\%$.
\begin{enumerate}[label=\alph*)]
  \item Calcule el precio actual $P_0$ suponiendo $g_2=4\%$.
  \item Repita el cálculo si el crecimiento perpetuo pasa a $g_2=5\%$ y comente la sensibilidad del valor.
  \item Con el precio del apartado (a), determine el rendimiento por dividendo y la ganancia de capital esperada durante el primer año y verifique que la suma coincide con la rentabilidad requerida.
\end{enumerate}
\end{ejercicio}

\section{Valoración de empresas y flujos de caja}

\begin{ejercicio}{4 (FCFF y FCFE: conciliación)}
Para una empresa se tiene: $\text{EBIT}=200$, $T_c=30\%$, $\text{Dep\&Amort}=25$, $\text{Capex}=60$, $\Delta\text{WorkingCap}=15$, $\text{NetIncome}=90$ y $\Delta\text{Deuda Neta}=10$ (todo en millones de euros).
\begin{enumerate}[label=\alph*)]
  \item Calcule $\text{FCFF}$
  \item Calcule $\text{FCFE}$
  % \item Verifique numéricamente que $\text{FCFF} - \text{Intereses netos después de impuestos} = \text{FCFE}$.
\end{enumerate}
\end{ejercicio}

\begin{ejercicio}{5 (Valoración por FCFF y WACC)}
Considere la empresa Pictor con $E=\,\EUR{600}$ millones, $D=\,\EUR{400}$ millones, $r_f=3\%$, $r_m=9\%$, $\beta=1.2$, $R_d=5\%$ y $T_c=30\%$. Los FCFF esperados (millones) son $\{95,110,125,140,150\}$, con crecimiento perpetuo de $3\%$ desde $t=5$. Determine:
\begin{enumerate}[label=\alph*)]
  \item El $\text{WACC}$.
  \item El valor de empresa descontando los FCFF y el valor terminal.
  % \item El valor del equity tras sumar caja excedente de \EUR{20} millones y restar la deuda.
  \item La misma empresa tiene $\text{FCFE}_1=85$, $\text{FCFE}_2=90$ y $\text{FCFE}_3=96$ (millones). Desde $t=3$, los flujos crecen al 2\% anual. El coste del equity es 10\% y existen 60 millones de acciones. Calcule el valor del equity hoy y obtenga el precio por acción
\end{enumerate}
\end{ejercicio}

\begin{ejercicio}{7 (Proyecto y sensibilidad al WACC)}
Una firma se financia con \EUR{140} millones de equity y \EUR{560} millones de deuda. $\beta=1.25$, $r_f=2.5\%$, $r_m-r_f=6\%$, $R_d=5.2\%$ y $T_c=30\%$. Se evalúa un proyecto con FCFF (millones) $\{55,60,68,75\}$, inversión inicial de \EUR{950} millones y crecimiento perpetuo del 2\% desde $t=4$.
\begin{enumerate}[label=\alph*)]
  \item Calcule el $\text{WACC}$.
  \item Obtenga el VAN del proyecto.
  \item Analice cómo cambiaría el VAN si el WACC aumentara o disminuyera en 1 punto porcentual.
\end{enumerate}
\end{ejercicio}

\section{Rentabilidad y riesgo individual}

\begin{ejercicio}{8 (Momentos y beta bajo escenarios)}
Considere los activos $A$ y $B$ con la siguiente distribución anual:
\begin{center}
\begin{tabular}{lccc}
\toprule
Escenario & Prob. & $r_A$ & $r_B$ \\
\midrule
Expansión & 0.30 & 18\% & 7\% \\
Base & 0.50 & 5\% & 4\% \\
Recesión & 0.20 & $-7\%$ & $-2\%$ \\
\bottomrule
\end{tabular}
\end{center}
Suponga que el rendimiento del mercado en cada escenario es $\{12\%, 6\%, -3\%\}$ respectivamente.
\begin{enumerate}[label=\alph*)]
  \item Calcule $\E[r_A]$, $\E[r_B]$, las varianzas y desviaciones estándar.
  \item Obtenga la covarianza y correlación entre $A$ y $B$.
  \item Calcule las betas $\beta_A$ y $\beta_B$ respecto al mercado.
\end{enumerate}
\end{ejercicio}

\begin{ejercicio}{9 (Rentabilidades históricas y ruido muestral)}
Una acción presenta la siguiente información mensual (en euros):
\begin{center}
\begin{tabular}{c c c c}
\toprule
Mes & Precio inicial & Precio final & Dividendo \\
\midrule
1 & 48.00 & 50.00 & 0.30 \\
2 & 50.00 & 49.00 & 0.30 \\
3 & 49.00 & 51.00 & 0.30 \\
4 & 51.00 & 53.00 & 0.30 \\
5 & 53.00 & 52.00 & 0.35 \\
6 & 52.00 & 55.00 & 0.35 \\
\bottomrule
\end{tabular}
\end{center}
\begin{enumerate}[label=\alph*)]
  \item Calcule las rentabilidades mensuales $R_t$.
  \item Obtenga la rentabilidad media muestral y la desviación estándar mensual y anualizada.
  % \item Comente la precisión de la estimación de la media con sólo seis observaciones.
\end{enumerate}
\end{ejercicio}

% \begin{ejercicio}{10 (CAPM, alfas y estrategia)}
% Con $r_f=3\%$ y $r_m=8.5\%$, tres acciones tienen $\beta_A=0.8$, $\beta_B=1.1$, $\beta_C=1.4$ y rentabilidades esperadas $\{7.2\%, 9.9\%, 12.1\%\}$.
% \begin{enumerate}[label=\alph*)]
%   \item Calcule la rentabilidad requerida por CAPM para cada activo.
%   \item Determine el \textit{alpha} y clasifique cada acción.
%   \item Proponga una estrategia de cartera de costo cero que explote los alfas positivos identificados.
% \end{enumerate}
% \end{ejercicio}

% \section{Carteras y diversificación}

% \begin{ejercicio}{11 (Cartera de dos activos y varianza mínima)}
% Para los activos del Ejercicio 8 ($\sigma_A^2=0.00582$, $\sigma_B^2=0.00050$, $\text{Cov}(A,B)=-0.00155$):
% \begin{enumerate}[label=\alph*)]
%   \item Calcule $\E[r_p]$ y $\sigma_p$ para $w_A=60\%$.
%   \item Determine el peso $w_A^*$ que minimiza la varianza y la desviación típica asociada.
% \end{enumerate}
% \end{ejercicio}

% \begin{ejercicio}{12 (Cartera de tres activos y ratio de Sharpe)}
% Considere tres activos con pesos $w = (0.40, 0.35, 0.25)$, rentabilidades esperadas $\{6\%, 8\%, 11\%\}$, desviaciones estándar $\{10\%, 14\%, 18\%\}$ y correlaciones $\rho_{12} = 0.45$, $\rho_{13} = 0.25$, $\rho_{23} = 0.50$. Sea $r_f = 2\%$.
% \begin{enumerate}[label=\alph*)]
%   \item Calcule la rentabilidad esperada y la desviación típica de la cartera.
%   \item Determine el ratio de Sharpe.
% \end{enumerate}
% \end{ejercicio}

% \begin{ejercicio}{13 (Correlación y diversificación)}
% Se forma una cartera equiponderada de dos activos con $\sigma_A = 18\%$ y $\sigma_B = 24\%$.
% \begin{enumerate}[label=\alph*)]
%   \item Calcule $\sigma_p$ para $\rho_{AB} = \{-0.3, 0, 0.5, 0.9\}$.
%   \item Repita el cálculo para pesos $w_A=0.6$, $w_B=0.4$ con $\rho_{AB} = 0.35$ y $-0.20$. Explique la diferencia respecto al caso equiponderado.
% \end{enumerate}
% \end{ejercicio}

% \begin{ejercicio}{14 (Límite de diversificación y riesgo sistemático)}
% Una familia de activos posee varianza individual $\sigma^2 = 0.04$ y covarianza media $\overline{\text{Cov}} = 0.01$. Para carteras equiponderadas de $N$ activos:
% \begin{enumerate}[label=\alph*)]
%   \item Calcule $\sigma_p$ para $N = 1, 5, 10, 30$.
%   \item Muestre que $\lim_{N\to\infty} \sigma_p = \sqrt{\overline{\text{Cov}}}$. Relacione el resultado con el concepto de riesgo sistemático (beta).
% \end{enumerate}
% \end{ejercicio}

\end{document}

